\chapter{Geschwindigkeitsqualifizierte Folgen}

Das laufende Alignment und sein pose3-Zustand beschreiben nun, wo das
Gleis-Frame liegt und wie es orientiert ist. Eine physische Antwort bestimmen
sie noch nicht. Geschwindigkeit tritt als getrennte Betriebseingabe hinzu;
Fahrzeugmodell und Bewertungsregel nochmals getrennt.

\section{Von der Bogenlänge zur Zeit}

Für konstantes \(v>0\) gilt
\[
\frac{\dd s}{\dd t}=v,
\qquad
\frac{\dd f}{\dd t}=v\frac{\dd f}{\dd s}.
\]
Damit erhält die räumliche Krümmungsrate \(\kappa'(s)\) erst nach Angabe der
Geschwindigkeit eine Zeitrate. Dasselbe gilt für Überhöhungsrate und höhere
Ableitungen.

\section{Eine reduzierte illustrative Beziehung}

Für die ebene Bewegung eines Punkts entlang der Referenztrajektorie gilt
\[
a_\kappa(s)=v^2\kappa(s).
\]
Unter dem lokalen rechtshändigen Frame, eng begrenzten
Starrkörperannahmen, der Überhöhungskonvention aus
\cref{eq:derived-roll-angle} und unter Vernachlässigung von Federung,
Rad--Schiene-Kontakt, Wind und weiteren Effekten lautet eine reduzierte
effektive Querkomponente
\[
a_{\mathrm{eff}}(s)=v^2\kappa(s)-g\sin\phi(s).
\]
Dies ist eine illustrative abgeleitete Größe, weder ein validiertes
vollständiges Fahrzeugdynamikmodell noch eine autoritative
Überhöhungsentwurfsregel.

\begin{table}[htbp]
\centering
\small
\begin{tabular}{lrr}
\toprule
Fall am selben \(s\) & \(v^2\kappa\) & \(a_{\mathrm{eff}}\) für \(\phi=4^\circ\)\\
\midrule
\(v=20\,\mathrm{m/s}\), \(\kappa=0.001\,\mathrm{m^{-1}}\) & \(0.40\,\mathrm{m/s^2}\) & \(-0.28\,\mathrm{m/s^2}\)\\
\(v=40\,\mathrm{m/s}\), \(\kappa=0.001\,\mathrm{m^{-1}}\) & \(1.60\,\mathrm{m/s^2}\) & \(0.92\,\mathrm{m/s^2}\)\\
\bottomrule
\end{tabular}
\caption{Illustrative Geschwindigkeitsabhängigkeit bei gleicher Krümmung und Überhöhung (\(g=9.81\,\mathrm{m/s^2}\), gerundet).}
\label{tab:same-curvature-two-speeds}
\end{table}

\Cref{tab:same-curvature-two-speeds} isoliert die Abhängigkeit: Weder Alignment
noch \(\kappa(s)\) ändern sich, doch eine Verdopplung der Geschwindigkeit
vervierfacht den krümmungsinduzierten Anteil. Eine Änderung der Überhöhung
würde effektive Komponente und Frame ändern, nicht das horizontale
Krümmungsgesetz.

\section{Änderung und betriebliche Bedeutung}

Unter denselben Annahmen und bei konstanter Geschwindigkeit gilt
\[
j(s)=v^3\kappa'(s)-gv\cos\phi(s)\,\phi'(s).
\]
Diese Beziehung erklärt, warum Übergangsbogenform und Überhöhungsrampe
betrieblich gemeinsam relevant sind und dennoch getrennte Gesetze bleiben.
Bei veränderlicher Geschwindigkeit treten zusätzliche Zeitableitungsterme auf.
Die Fahrzeugantwort hängt weiter von Masse, Federung, Rad--Schiene-Kontakt und
anderen Modellentscheidungen ab. Keine dieser Größen wird stillschweigend in
die Alignment-Geometrie aufgenommen.

\section{Bewertung ist eine spätere Operation}

Die Abhängigkeitsreihenfolge lautet
\begin{align*}
\text{Geometrie}&\rightarrow\text{realisiertes Frame}
\rightarrow\text{Geschwindigkeitsprofil}\rightarrow\text{Fahrzeugmodell}\\
&\rightarrow\text{abgeleitete Dynamik}
\rightarrow\text{Zweck und Regel}\rightarrow\text{Bewertung}.
\end{align*}
Ein Grenzwert aus einer Norm ist nur unter deren Anwendungsbereich und Version
anwendbar. Eine berechnete Überschreitung ist ein Bewertungsergebnis; Annahme,
Ablehnung oder Autorisierung benötigen eine Engineering Decision mit der
zuständigen Autorität.

\begin{summary}
Geschwindigkeit ändert die Folge des laufenden Alignments, nicht seine
intrinsische Geometrie. Überhöhung ändert Gleis-Frame und reduzierte effektive
Wirkung, während Geometrie, Betriebseingabe, Fahrzeugmodell, Ableitung,
Bewertung und Entscheidung getrennt bleiben.
\end{summary}
